
\documentclass[tikz, border=10mm]{standalone}
\usepackage{tikz}
\usetikzlibrary{positioning, calc, backgrounds}
% nn_blocks.tex
% Defines styles for neural network blocks using TikZ
% Now supports Hybrid 2D/3D Visualization with PlotNeuralNet

% --- PlotNeuralNet Integration ---
% We assume Ball.sty, Box.sty, RightBandedBox.sty are in the same directory
\usepackage{Ball}
\usepackage{Box}
\usepackage{RightBandedBox}

% PlotNeuralNet Init Definitions (from init.tex)
\usetikzlibrary{quotes,arrows.meta}
\usetikzlibrary{positioning}
\def\edgecolor{rgb:blue,4;red,1;green,4;black,3}
\newcommand{\midarrow}{\tikz \draw[-Stealth,line width =0.8mm,draw=\edgecolor] (-0.3,0) -- ++(0.3,0);}

% Define PlotNeuralNet Colors (Standard Palette)
\definecolor{ConvColor}{rgb}{0.98,0.81,0.69}  % Apricot/Orange
\definecolor{ConvReluColor}{rgb}{0.98,0.81,0.69}
\definecolor{PoolColor}{rgb}{0.8,0.33,0.33}   % Redish
\definecolor{UnpoolColor}{rgb}{0.6,0.7,0.9}   % Blueish
\definecolor{FcColor}{rgb}{0.6,0.2,0.8}       % Purple
\definecolor{FcReluColor}{rgb}{0.6,0.2,0.8}
\definecolor{SoftmaxColor}{rgb}{0.9,0.9,0.6}  % Yellowish
\definecolor{SumColor}{rgb}{0.9,0.9,0.9}      % Gray

% Custom colors for our specific blocks
\definecolor{ResBlockColor}{rgb}{0.85,0.92,1.0} % Light Blue
\definecolor{EmbedColor}{rgb}{0.95,0.95,0.8}    % Light Yellow
\definecolor{TransColor}{rgb}{0.8,0.9,0.8}      % Light Green

% --- NNTikZ 2D Styles (Matched to EDSR Style) ---
\usetikzlibrary{positioning, chains, shapes.geometric, fit, shapes, arrows.meta, calc, backgrounds, shadows}

\tikzset{
    >=LaTeX, 
    very thick,
    node distance=1.0cm and 1.2cm,
    % Basic Arrow
    arrow/.style={
        ->,
        very thick,
        draw=black!80,
        rounded corners=0.2cm
    },
    % Dashed Skip Connection
    skip/.style={
        arrow,
        dashed,
        draw=blue!60
    },
    % Grouping Block (Background Box)
    block/.style={
        rectangle,
        fill=gray!10,
        rounded corners=3mm,
        draw=black!60,
        very thick,
        inner sep=0.8em,
        drop shadow={opacity=0.15}
    },
    nntikz_block/.style={block}, % Alias
    % General Layer
    layer/.style={
        rectangle,
        fill=white!10,
        rounded corners=1mm,
        inner xsep=0.6em,
        inner ysep=0.6em,
        minimum height=3.0em,
        align=center,
        draw=black!80,
        very thick,
        drop shadow={opacity=0.2, shadow xshift=1mm, shadow yshift=-1mm}
    },
    nntikz_layer/.style={layer}, % Alias
    % Specific Layer Types
    conv/.style={
        layer,
        fill=orange!10,
        draw=orange!80!black
    },
    relu/.style={
        layer,
        fill=yellow!10,
        draw=yellow!80!black,
        minimum height=2.5em
    },
    pool/.style={
        layer,
        fill=red!10,
        draw=red!80!black
    },
    up/.style={
        layer,
        fill=purple!10,
        draw=purple!80!black
    },
    res/.style={ % Residual Block or similar
        layer,
        fill=blue!5,
        draw=blue!80!black
    },
    op/.style={ % Other operations
        layer,
        fill=green!10,
        draw=green!80!black
    },
    % Input/Output
    input/.style={ 
        circle,
        minimum width=3.0em,
        draw=black!80,
        fill=gray!20,
        thick,
        align=center,
        drop shadow={opacity=0.2}
    },
    io/.style={input}, % Alias
    % Summation
    sum/.style={
        circle,
        draw=black!80,
        fill=white,
        inner sep=0pt,
        minimum size=1.2em,
        node contents={+}
    }
}



\begin{document}
\begin{tikzpicture}[node distance=1.5cm, every node/.style={transform shape}]

% Title
\node[font=\Large\bfseries] (title) at (5, 8) {Architecture Legend};

% --- 3D Section ---
\node[font=\bfseries, right] (sec3d) at (0, 7) {3D Visualization (PlotNeuralNet)};

% 3D Elements (Manually positioning Boxes)
% Note: PlotNeuralNet uses a shifted coordinate system. We'll use simple TikZ nodes with 3D styling for the legend
% to avoid complex projection issues, or use \pic if possible.

% Conv
\coordinate (l3d_start) at (1, 6);
\pic[shift={(0,0,0)}] at (l3d_start) {Box={name=l_conv, fill=\ConvColor, height=3, width=1, depth=3, caption= }};
\node[right=0.5cm of l_conv-east, anchor=west] {Conv / ConvReLU};

% Pool
\coordinate (l_pool_pos) at ($(l3d_start) + (5, 0)$);
\pic[shift={(0,0,0)}] at (l_pool_pos) {Box={name=l_pool, fill=\PoolColor, height=2, width=1, depth=2, caption= }};
\node[right=0.5cm of l_pool-east, anchor=west] {Pooling};

% Up
\coordinate (l_up_pos) at ($(l_pool_pos) + (5, 0)$);
\pic[shift={(0,0,0)}] at (l_up_pos) {Box={name=l_up, fill=\UnpoolColor, height=3, width=1, depth=3, caption= }};
\node[right=0.5cm of l_up-east, anchor=west] {Upsample / Unpool};

% Next Row 3D
\coordinate (l3d_row2) at (1, 3.5);

% FC / Transformer
\pic[shift={(0,0,0)}] at (l3d_row2) {Box={name=l_fc, fill=\FcColor, height=1, width=4, depth=1, caption= }};
\node[right=0.5cm of l_fc-east, anchor=west] {Linear / Transformer};

% ResBlock
\coordinate (l_res_pos) at ($(l3d_row2) + (6, 0)$);
\pic[shift={(0,0,0)}] at (l_res_pos) {Box={name=l_res, fill=\ResBlockColor, height=2.5, width=1.5, depth=2.5, caption= }};
\node[right=0.5cm of l_res-east, anchor=west] {Residual Block};

% --- 2D Section ---
\node[font=\bfseries, right] (sec2d) at (0, 1) {2D Visualization (Flowchart)};

\coordinate (l2d_start) at (1, 0);

% Conv
\node[conv, minimum width=1.5cm] (n_conv) at (l2d_start) {Conv};
\node[right=0.2cm of n_conv] {Convolution};

% ReLU
\node[relu, minimum width=1.5cm, right=3cm of n_conv] (n_relu) {Act};
\node[right=0.2cm of n_relu] {Activation (ReLU/GELU)};

% Pool
\node[pool, minimum width=1.5cm, right=3.5cm of n_relu] (n_pool) {Pool};
\node[right=0.2cm of n_pool] {Pooling};

% Next Row 2D
\coordinate (l2d_row2) at (1, -2);

% Up
\node[up, minimum width=1.5cm] (n_up) at (l2d_row2) {Up};
\node[right=0.2cm of n_up] {Upsample};

% Res
\node[res, minimum width=1.5cm, right=3cm of n_up] (n_res) {Res};
\node[right=0.2cm of n_res] {Residual/Swin Block};

% Op
\node[op, minimum width=1.5cm, right=3.5cm of n_res] (n_op) {Op};
\node[right=0.2cm of n_op] {Other Op};

% Skip Connection
\coordinate (skip_start) at (1, -4);
\coordinate (skip_end) at (3, -4);
\draw[skip] (skip_start) -- (skip_end) node[midway, above] {Skip Connection};
\node[right=0.2cm of skip_end] {Residual / Skip Connection};

\end{tikzpicture}
\end{document}
